\JOUChapterWithEnglish{分选过程响应与优化}{Separation Response and Optimization}
\label{chap:optimization}

\JOUSectionWithEnglish{操作参数耦合影响}{Coupled Effects of Operating Parameters}

在实验尺度和工程尺度下，循环矿浆量、给矿浓度、充气量和背压并不是彼此独立的单因素，而是通过流场组织、气泡更新频率和颗粒停留时间形成耦合响应。采用响应面分析与参数扫描相结合的方法，可以更有效地识别不同工况下的优势操作窗口。

\vspace{0.4\baselineskip}

优化过程的核心不在于单纯追求某一指标的极值，而在于协调精矿品位、回收率和运行稳定性之间的平衡。将样机数据与数值模拟结果进行联合校核，有助于提升模型参数的可迁移性和工艺放大时的鲁棒性。
